% Options for packages loaded elsewhere
\PassOptionsToPackage{unicode}{hyperref}
\PassOptionsToPackage{hyphens}{url}
\documentclass[
  journal]{IEEEtran}
\usepackage{xcolor}
\usepackage{amsmath,amssymb}
\setcounter{secnumdepth}{-\maxdimen} % remove section numbering
\usepackage{iftex}
\ifPDFTeX
  \usepackage[T1]{fontenc}
  \usepackage[utf8]{inputenc}
  \usepackage{textcomp} % provide euro and other symbols
\else % if luatex or xetex
  \usepackage{unicode-math} % this also loads fontspec
  \defaultfontfeatures{Scale=MatchLowercase}
  \defaultfontfeatures[\rmfamily]{Ligatures=TeX,Scale=1}
\fi
\usepackage{lmodern}
\ifPDFTeX\else
  % xetex/luatex font selection
\fi
% Use upquote if available, for straight quotes in verbatim environments
\IfFileExists{upquote.sty}{\usepackage{upquote}}{}
\IfFileExists{microtype.sty}{% use microtype if available
  \usepackage[]{microtype}
  \UseMicrotypeSet[protrusion]{basicmath} % disable protrusion for tt fonts
}{}
\makeatletter
\@ifundefined{KOMAClassName}{% if non-KOMA class
  \IfFileExists{parskip.sty}{%
    \usepackage{parskip}
  }{% else
    \setlength{\parindent}{0pt}
    \setlength{\parskip}{6pt plus 2pt minus 1pt}}
}{% if KOMA class
  \KOMAoptions{parskip=half}}
\makeatother
\setlength{\emergencystretch}{3em} % prevent overfull lines
\providecommand{\tightlist}{%
  \setlength{\itemsep}{0pt}\setlength{\parskip}{0pt}}
\usepackage{bookmark}
\IfFileExists{xurl.sty}{\usepackage{xurl}}{} % add URL line breaks if available
\urlstyle{same}
\hypersetup{
  pdftitle={Locality Dominance in Hyperspectral and Multimodal Land-Cover Classification},
  pdfauthor={Will Dale; Pond Science Institute},
  hidelinks,
  pdfcreator={LaTeX via pandoc}}

\title{Locality Dominance in Hyperspectral and Multimodal Land-Cover
Classification}
\author{Will Dale}
\date{}

\begin{document}
\maketitle

\begin{abstract}

Hyperspectral and multimodal land-cover classification methods typically
rely on high-dimensional per-pixel spectral representations and
increasingly complex fusion strategies. In this work, we demonstrate
that a compact neighborhood-statistic representation---computed entirely
from local spatial context---outperforms substantially
higher-dimensional spectral baselines across three benchmarks. Using
stratified 10\% training splits, a 10--14 dimensional patch-statistics
feature vector (mean, standard deviation, gradient magnitude, and
PCA-derived neighborhood descriptors) achieves +3.0--11.1 percentage
point (pp) improvements in overall accuracy (OA) and +6.3--16.1pp in
average accuracy (AA) relative to 30--200 dimensional spectral baselines
on Houston (multimodal), Indian Pines, and PaviaU datasets. Notably,
gains in AA exceed gains in OA, indicating improved minority-class
performance. Accuracy saturates at a small, scene-dependent neighborhood
scale (5\(\times\)5--7\(\times\)7), beyond which larger windows offer no consistent
benefit. In contrast, uncertainty-driven gating and spectral transform
augmentations provide negligible or negative impact once spatial context
is encoded. These results suggest that local second-order spatial
statistics constitute a sufficient and dominant representation for class
separation in standard hyperspectral benchmarks, substantially reducing
feature dimensionality while improving performance.

\end{abstract}

\section{Introduction}

Hyperspectral image (HSI) and multimodal remote sensing classification
remain central tasks in land-cover analysis. Traditional approaches
emphasize high-dimensional spectral representations, often employing
dimensionality reduction (e.g., PCA), spectral transforms, or
increasingly complex multimodal fusion strategies. More recently, deep
convolutional architectures have demonstrated strong performance by
implicitly exploiting spatial context through learned receptive fields.
Despite these advances, the relative importance of spectral transforms
versus explicit spatial aggregation remains underexplored.

Most per-pixel spectral classifiers operate on high-dimensional feature
vectors (100--200 bands), implicitly assuming that class separability is
encoded primarily in spectral signatures. However, real-world scenes
exhibit substantial intra-class spectral variability arising from
illumination changes, mixed pixels, and local texture effects.
Consequently, pixel-wise models frequently require either high model
capacity or complex fusion mechanisms to stabilize predictions.

In this work, we revisit a fundamental question: \textbf{Is
high-dimensional spectral encoding necessary once local spatial context
is explicitly represented?}

We evaluate a compact neighborhood-statistics representation constructed
entirely from local spatial summary features---mean, standard deviation,
gradient magnitude, and PCA-derived descriptors computed within small
spatial windows. The resulting feature vector is 10--14 dimensions,
independent of the number of spectral bands.

Across three benchmarks---Houston (multimodal HSI+MS+LiDAR), Indian
Pines, and PaviaU---using stratified 10\% training splits, we observe
that this neighborhood-based representation consistently outperforms
30--200 dimensional spectral baselines. Improvements range from +3.0 to
+11.1pp in overall accuracy and +6.3 to +16.1pp in average accuracy,
with substantial gains in minority-class performance. Accuracy saturates
at a small, scene-dependent window size (5\(\times\)5--7\(\times\)7), indicating the
existence of a minimal sufficient spatial scale.

We further show that uncertainty-driven gating strategies do not improve
performance in the multimodal setting and that spectral transform
augmentations become redundant once spatial context is encoded.

The primary contributions of this work are:

\begin{enumerate}
\def\labelenumi{\arabic{enumi}.}
\tightlist
\item
  \textbf{Neighborhood Sufficiency Principle (empirical):} A compact
  local-statistics representation achieves performance comparable to or
  exceeding high-dimensional spectral encodings across datasets.
\item
  \textbf{Minimal Spatial Scale Identification:} Classification accuracy
  saturates at a small, scene-dependent neighborhood radius r*, ranging
  from 5\(\times\)5 to 7\(\times\)7 across tested benchmarks.
\item
  \textbf{Minority-Class Stabilization:} Neighborhood statistics
  disproportionately improve average accuracy (\(\Delta\)AA \textgreater{} \(\Delta\)OA),
  mitigating majority-class bias inherent in high-dimensional spectral
  models.
\item
  \textbf{Negative Fusion Result:} Uncertainty-based gating (entropy,
  margin, Gini) does not confer benefit in the evaluated multimodal
  benchmark; gate confidence exhibits slight negative correlation with
  classification correctness.
\end{enumerate}

These findings suggest that local second-order spatial statistics
constitute a dominant representational mechanism for standard
hyperspectral benchmarks, with significant dimensionality and
computational advantages.

\section{Related Work}\label{related-work}

\subsection{Pixel-wise Spectral
Classification}\label{pixel-wise-spectral-classification}

Classical hyperspectral image (HSI) classification methods operate on
per-pixel spectral signatures, often comprising 100--200 bands.
Dimensionality reduction techniques such as principal component analysis
(PCA), linear discriminant analysis (LDA), and band selection methods
are commonly applied prior to classification using support vector
machines (SVM), random forests (RF), or k-nearest neighbors. While
effective in controlled settings, pixel-wise approaches are sensitive to
intra-class spectral variability and mixed pixels, often requiring high
model capacity to compensate.

Our spectral baselines (PCA-30 and raw spectral RF) follow this standard
protocol to provide a direct comparison against neighborhood-based
representations.

\subsection{Spectral--Spatial Feature
Extraction}\label{spectralspatial-feature-extraction}

To address the limitations of pixel-only models, numerous works
incorporate spatial context through handcrafted texture descriptors,
morphological profiles, and spatial filtering techniques. Extended
morphological profiles (EMP), gray-level co-occurrence matrix (GLCM)
features, and local binary patterns (LBP) have demonstrated that
second-order spatial statistics improve classification stability.

More recently, patch-based feature extraction and joint
spectral--spatial descriptors have become standard, particularly in
conjunction with deep architectures. However, many methods either
dramatically increase feature dimensionality or rely on learned
convolutional filters.

In contrast, the present work isolates a minimal set of low-dimensional
neighborhood summary statistics and evaluates their sufficiency without
deep models or high-capacity transforms.

\subsection{Convolutional Neural Networks for
HSI}\label{convolutional-neural-networks-for-hsi}

Convolutional neural networks (CNNs) have achieved state-of-the-art
performance in hyperspectral classification by learning hierarchical
spectral--spatial filters. These architectures implicitly exploit local
receptive fields, progressively aggregating spatial information across
layers.

While CNNs demonstrate strong performance, they introduce significant
computational overhead and model complexity. The present study
investigates whether the performance gains attributed to deep models may
arise primarily from local spatial aggregation rather than from deep
spectral transformation capacity.

Our results suggest that small fixed neighborhoods (5\(\times\)5--7\(\times\)7) with
simple second-order statistics capture much of the discriminative power
typically associated with convolutional architectures on standard
benchmarks.

\subsection{Multimodal Fusion and
Gating}\label{multimodal-fusion-and-gating}

Multimodal remote sensing benchmarks (e.g., HSI + LiDAR + multispectral)
have motivated diverse fusion strategies, including feature
concatenation, late fusion, attention mechanisms, and uncertainty-driven
gating. Entropy-based or margin-based gating schemes aim to weight
modalities according to classifier confidence.

However, uncertainty estimates may not align with true class-separating
invariants in complex scenes. Our results provide empirical evidence
that uncertainty-driven gating can degrade performance and that
amplitude-based heuristics provide only modest improvement relative to
explicit spatial encoding.

\section{Methods}\label{methods}

\subsection{Datasets}\label{datasets}

We evaluate on three benchmarks:

\begin{itemize}
\tightlist
\item
  \textbf{Houston Multimodal:} Hyperspectral (144 bands), multispectral
  (8 bands), and LiDAR elevation data with 15 classes. Pre-extracted
  11\(\times\)11 patches are provided; 2832 labeled samples total.
\item
  \textbf{Indian Pines:} 145\(\times\)145 HSI with 200 spectral bands and 16
  classes (10249 labeled pixels).
\item
  \textbf{PaviaU:} 610\(\times\)340 HSI with 103 spectral bands and 9 classes
  (42776 labeled pixels).
\end{itemize}

For Indian Pines and PaviaU, patches are extracted dynamically from the
full image cube using reflect padding at image borders.

\subsection{Train--Test Protocol}\label{traintest-protocol}

For Indian Pines and PaviaU, we use stratified random splits with 10\%
of labeled pixels per class allocated to training and 90\% to testing,
following standard literature protocol. For Houston multimodal, we use a
70\%/30\% stratified split. All reported results use identical splits
across methods for fair comparison. For Houston and Indian Pines,
experiments are repeated across three random seeds (42, 1, 7) to assess
stability.

PCA transformations are fitted exclusively on the training set and
applied to test data to prevent leakage. Patch-level PCA statistics are
likewise fitted on training patch pixels only.

\subsection{Spectral Baselines}\label{spectral-baselines}

Two spectral baselines are evaluated:

\begin{enumerate}
\def\labelenumi{\arabic{enumi}.}
\tightlist
\item
  \textbf{PCA-k spectral baseline:} PCA reduced to k components (k=30
  for HSI-only datasets, k=50 for Houston), classified using Random
  Forest (200 estimators).
\item
  \textbf{Raw spectral baseline:} All bands (103--200) used directly
  with Random Forest.
\end{enumerate}

For Houston multimodal, a 59-dimensional concatenation baseline (HSI
PCA-50 + LiDAR + MS) is used as the primary multimodal reference, along
with area-weighted gating variants.

\subsection{Neighborhood-Statistic
Representation}\label{neighborhood-statistic-representation}

For each labeled pixel, a square window of size r\(\times\)r centered at the
pixel is extracted. The following features are computed over the
neighborhood:

\begin{itemize}
\tightlist
\item
  Mean over all spectral bands (flattened patch)
\item
  Standard deviation over all spectral bands
\item
  Gradient magnitude (finite differences on mean-pooled spectral
  intensity at patch center)
\item
  Mean and standard deviation of the first three PCA components computed
  over the r\(\times\)r patch pixels (6 features)
\end{itemize}

For the multimodal Houston dataset, additional modality-specific
statistics are included: LiDAR center-pixel residual (center minus patch
mean), MS gradient magnitude, and per-modality mean/std. This yields a
14-dimensional feature vector for Houston and a 10-dimensional vector
for HSI-only datasets.

Window sizes r \(\in\) \{3, 5, 7, 11\} are evaluated to identify the minimal
sufficient spatial scale. The chosen center for all crops is pixel
position (5, 5) within the 11\(\times\)11 pre-extracted patch for Houston, and
the labeled pixel center for dynamically extracted patches.

\subsection{Uncertainty Gating (Houston
Only)}\label{uncertainty-gating-houston-only}

For the Houston multimodal dataset, per-modality Random Forest
classifiers are trained and used to compute uncertainty signals per
sample: entropy (H = -\(\sum\) p log p), margin (top-1 minus top-2
probability), and Gini impurity (1 - \(\sum p^2\)). Gate weights are computed as
softmax(u/T) where u is the uncertainty signal and T is a temperature
parameter. Sweeps over T \(\in\) \{0.1, 0.5, 1.0, 2.0, 5.0\} are conducted.

Gate diagnostics reported include: per-modality mean gate weight,
saturation rate (fraction of samples where any gate weight
\textgreater{} 0.9 or \textless{} 0.1), and Pearson correlation between
maximum gate weight and per-sample classification correctness.

\subsection{Classification Model}\label{classification-model}

Random Forest classifiers (200 estimators, fixed random seed) are used
for all primary comparisons to isolate representational effects from
model capacity. Experiments with gradient-boosted decision trees
(HistGradientBoostingClassifier) confirm that the ranking between
representations is stable across classifier choices (difference \(\leq\)
0.24pp).

\subsection{Evaluation Metrics}\label{evaluation-metrics}

Performance is measured using:

\begin{itemize}
\tightlist
\item
  \textbf{Overall Accuracy (OA):} Fraction of correctly classified test
  pixels.
\item
  \textbf{Average Accuracy (AA):} Mean of per-class accuracies;
  emphasizes minority-class performance under class imbalance.
\item
  \textbf{Cohen's kappa coefficient (\(\kappa\)):} Agreement corrected for
  chance.
\end{itemize}

Average Accuracy is emphasized throughout to assess whether gains are
broad-based or driven by dominant classes.

\section{Results}\label{results}

\subsection{Spectral Baselines and Gating Methods (Houston
Multimodal)}\label{spectral-baselines-and-gating-methods-houston-multimodal}

Uncertainty-driven gating degrades classification performance on the
Houston multimodal benchmark. Entropy gating (T = 1.0) reduces overall
accuracy (OA) by 9.11 percentage points (pp) relative to the
59-dimensional concatenation baseline (OA = 96.34\%), and margin gating
at T = 5.0 yields a further -0.48pp. The failure mechanism is
structural: the correlation between gate confidence (maximum modality
weight) and per-sample classification correctness is corr(max\_w,
correct) \(\approx\) -0.08, indicating that uncertainty is not a valid fusion
invariant on this manifold. The gate fires hardest precisely where the
classifier is most likely to be wrong, producing category-level damage
concentrated in spectrally ambiguous classes 10--13, which suffer
declines of 14--31pp relative to the unweighted baseline. Area gating,
which uses amplitude as a proxy for signal stability rather than
classifier uncertainty, is the sole gating variant to improve over
concat, providing a modest +0.35pp gain (OA = 96.69\%). These results
establish that naive uncertainty quantification does not confer fusion
benefit in this multimodal scene configuration.

\subsection{Neighborhood-Statistic
Representation}\label{neighborhood-statistic-representation-1}

Rather than aggregating modality-specific spectral transforms, the
patch-statistics representation encodes local spatial context directly
from the hyperspectral cube. Each sample is described by a
10--14-dimensional feature vector comprising patch mean, standard
deviation, gradient magnitude, and the leading PCA components computed
within a square spatial window centered on the labeled pixel. This
representation is constructed without any center-pixel spectral
information; the entire feature budget is allocated to neighborhood
summary statistics.

On the Houston multimodal benchmark, the patch-only configuration (14D,
5\(\times\)5 window) achieves a mean OA of 99.37\% across three independent
random seeds (seed = 42, 1, 7), compared to 96.34\% for the
59-dimensional concatenation baseline, a gain of +3.03pp (+2.7\%
relative improvement). Full per-seed and per-method results are reported
in Table 1. Augmenting the patch feature vector with chromogeometry
descriptors and amplitude anchors (27D total) reduces mean OA by 0.47pp
relative to patch-only, confirming that DFT-derived quadrances and
spectral centroids are redundant once local spatial statistics are
present. Substituting a gradient-boosted decision tree (GBDT) classifier
for the random forest adds at most 0.24pp and does not change the
ranking of feature configurations. The full patch-only pipeline requires
0.45 seconds of training time on this dataset, reflecting the low
dimensionality of the representation.

\subsection{Optimal Neighborhood
Scale}\label{optimal-neighborhood-scale}

A patch-size sweep on the Houston scene identifies r* = 5\(\times\)5 as the
optimal window. At 3\(\times\)3, patch-only OA is 97.64\%, a useful improvement
over the 96.34\% concat baseline but 2.24pp below the 5\(\times\)5 result,
indicating that a three-pixel window is insufficient to capture the
spatial homogeneity structure of this scene. Performance at 5\(\times\)5 (OA =
99.88\%, seed = 42) and 11\(\times\)11 (OA = 99.17\%) are within the range of
seed-to-seed variance, placing the plateau at r \(\geq\) 5 on this benchmark.
The absence of further gain beyond 5\(\times\)5 is consistent with class-boundary
contamination: as the window expands, mixed-class neighborhoods
introduce spectral mixture bias that offsets the variance reduction
benefit of averaging over a larger sample of within-class pixels.

\subsection{Cross-Dataset Generalization (Indian Pines and
PaviaU)}\label{cross-dataset-generalization-indian-pines-and-paviau}

The generalization experiment, conducted on Indian Pines (145\(\times\)145
pixels, 200 bands, 16 classes, 10\% training) and PaviaU (610\(\times\)340
pixels, 103 bands, 9 classes, 10\% training), constitutes the primary
evidence for locality dominance as a scene-independent phenomenon.

On Indian Pines, a 30-dimensional PCA spectral representation applied to
the center pixel achieves OA = 70.66\%, AA = 56.39\%, \(\kappa\) = 0.657 as the
reference spectral baseline. The 10-dimensional patch-only
representation at 7\(\times\)7 achieves OA = 81.00\%, AA = 72.46\%, \(\kappa\) = 0.782
(seed = 42), with a mean OA of 81.05\% across seeds 42, 1, and 7,
confirming seed stability. The gain in average accuracy (\(\Delta\)AA = +16.08pp)
substantially exceeds the gain in overall accuracy (\(\Delta\)OA = +10.34pp),
establishing that improvement is broad-based across classes and
disproportionately benefits rare and spectrally ambiguous categories.
This asymmetry is diagnostic: the 200-dimensional raw spectral baseline
achieves OA = 75.02\% but only AA = 60.18\%, reflecting majority-class
dominance that the spatial representation resolves. On PaviaU, the same
patch-only configuration (10D, 7\(\times\)7) yields OA = 91.98\%, AA = 87.96\%, \(\kappa\)
= 0.893, compared to the PCA-30 baseline of OA = 88.07\%, AA = 81.61\%,
\(\kappa\) = 0.837, a gain of \(\Delta\)OA = +3.91pp and \(\Delta\)AA = +6.35pp. Full OA, AA, and \(\kappa\)
results for both datasets across all patch sizes are reported in Table 2
and Table 3.

The optimal window on both hyperspectral-only benchmarks is r* = 7\(\times\)7,
one step larger than the r* = 5\(\times\)5 optimum on the Houston multimodal
scene. This shift is consistent with the larger spatial homogeneity
scale expected in agricultural landscapes (Indian Pines) and structured
urban blocks (PaviaU) relative to the densely textured Houston scene,
where the multimodal LiDAR channel provides additional structural
disambiguation.

\subsection{Summary}\label{summary}

Across all three datasets and a total of nine seed-configuration
combinations, a 10--14-dimensional neighborhood-statistic representation
consistently and substantially outperforms spectral baselines ranging
from 30 to 200 dimensions, achieving a feature-dimension reduction of
76--95\% while improving OA by 3.0--11.1pp and AA by 6.3--16.1pp. The
neighborhood operator is the dominant contributor to classification
performance; spectral transforms, uncertainty-based modality gating, and
chromogeometry descriptors are each redundant or mildly harmful once
spatial context is encoded. The optimal neighborhood scale r* is
scene-dependent and tracks the spatial homogeneity scale of the
land-cover pattern, saturating at 5\(\times\)5 for dense multimodal urban imagery
and at 7\(\times\)7 for hyperspectral agricultural and urban-block scenes.

\section{Discussion}\label{discussion}

\subsection{Why Local Neighborhood Statistics Outperform Spectral
Transforms}\label{why-local-neighborhood-statistics-outperform-spectral-transforms}

The central empirical finding across all three datasets is that
low-dimensional neighborhood statistics outperform substantially
higher-dimensional per-pixel spectral representations. This behavior can
be understood through variance reduction and spatial stationarity.

Let a pixel spectrum be modeled as:

\[x_i = \mu_c + \varepsilon_i\]

where \(\mu_c\) is the class-conditional mean and \(\varepsilon_i\) captures illumination
variability, sensor noise, subpixel mixing, and fine-scale texture.
Pixel-wise classifiers must separate classes under high within-class
variance.

By contrast, the neighborhood operator computes summary statistics over
an r\(\times\)r window:

\[\bar{x}_i = \frac{1}{|N_r|} \sum_{j \in N_r(i)} x_j\]

Under spatial coherence assumptions, variance shrinks approximately as
1/\textbar N\_r\textbar, while the class mean remains stable. The
representation therefore increases signal-to-noise ratio without
increasing feature dimensionality.

The PCA-derived patch statistics and local variance features
additionally encode second-order texture structure, capturing
discriminative information that spectral basis transforms (e.g.,
DFT-derived quadrances) do not uniquely provide once spatial context is
incorporated.

\subsection{The Neighborhood Sufficiency Principle
(Empirical)}\label{the-neighborhood-sufficiency-principle-empirical}

Across datasets, performance improves monotonically with window size up
to a scene-dependent threshold r*, beyond which accuracy plateaus or
slightly regresses.

This behavior is consistent with two competing effects:

\begin{enumerate}
\def\labelenumi{\arabic{enumi}.}
\tightlist
\item
  \textbf{Variance reduction} within homogeneous class regions.
\item
  \textbf{Boundary contamination} when windows span multiple classes.
\end{enumerate}

We therefore observe a minimal sufficient neighborhood scale:

\[r^* \approx \text{class spatial coherence radius}\]

Empirically: - Houston multimodal: r* \(\approx\) 5\(\times\)5 - Indian Pines: r* \(\approx\) 7\(\times\)7 -
PaviaU: r* \(\approx\) 7\(\times\)7

These differences align with known scene structure: agricultural fields
and urban blocks exhibit larger spatial homogeneity than dense mixed
scenes.

\textbf{Neighborhood Sufficiency Principle (empirical):}

\begin{quote}
A compact local-statistics representation computed over a small,
scene-dependent neighborhood achieves performance comparable to or
exceeding high-dimensional per-pixel spectral representations, with
performance saturating beyond a bounded spatial scale r*.
\end{quote}

\subsection{Minority-Class Rescue and Class
Imbalance}\label{minority-class-rescue-and-class-imbalance}

On Indian Pines, the improvement in average accuracy (AA) substantially
exceeds the improvement in overall accuracy (OA). This indicates that
gains are not limited to dominant classes but instead reflect broad
improvements across rare and spectrally ambiguous categories.

Pixel-only spectral models exhibit majority-class bias: OA increases
while AA remains suppressed. Neighborhood statistics mitigate this
effect by stabilizing minority-class signatures through local averaging
and texture encoding.

This suggests that spatial context functions not only as a
noise-reduction operator but also as a class-balance equalizer.

\subsection{Why Uncertainty Gating
Fails}\label{why-uncertainty-gating-fails}

Uncertainty-based gating assumes that per-modality classifier confidence
is aligned with classification correctness. However, in the Houston
multimodal dataset, gate confidence exhibits a slight negative
correlation with correctness (corr \(\approx\) -0.08). This indicates that
ambiguity is intrinsic to the class manifold rather than a reliable
fusion signal.

Area-based gating succeeds modestly because amplitude proxies local
signal stability. However, even this approach is dominated by explicit
spatial encoding.

Thus, uncertainty is not a reliable invariant for multimodal fusion in
these scenes; spatial locality is.

\subsection{Relation to Deep Learning and Texture
Theory}\label{relation-to-deep-learning-and-texture-theory}

The empirical dominance of small neighborhood statistics aligns with
classical texture descriptors (e.g., second-order statistics) and with
convolutional neural network receptive field theory.

CNNs implicitly learn spatial filters that capture local mean, variance,
and gradient information over small receptive fields. The present
results show that much of this benefit can be achieved explicitly using
handcrafted low-dimensional statistics, without deep architectures.

This suggests that the primary structural advantage of CNN-based
hyperspectral models may lie in their implicit spatial aggregation
rather than in high-capacity spectral transforms.

\subsection{Limitations and Scope}\label{limitations-and-scope}

The findings are empirical and evaluated on three benchmarks. While
locality dominance is consistent across datasets tested, the optimal
window size varies with scene structure. Further evaluation on
additional multimodal benchmarks and scenes with highly fragmented
spatial patterns would strengthen generality claims.

Additionally, the current work uses fixed window shapes and summary
statistics; adaptive or anisotropic neighborhoods may further improve
performance.

\subsection{Implications}\label{implications}

These results suggest a design principle for hyperspectral and
multimodal classification:

\begin{enumerate}
\def\labelenumi{\arabic{enumi}.}
\tightlist
\item
  Encode spatial context first.
\item
  Keep representation compact.
\item
  Introduce spectral transforms only if they add orthogonal information
  beyond local statistics.
\end{enumerate}

In the tested benchmarks, the neighborhood operator subsumes spectral
transform generators in class-separating power.

\section{Conclusion}\label{conclusion}

We investigated whether high-dimensional spectral representations are
necessary for hyperspectral and multimodal land-cover classification
once local spatial context is explicitly encoded. Across three
benchmarks (Houston multimodal, Indian Pines, and PaviaU) using
stratified 10\% training splits, we found that a compact 10--14
dimensional neighborhood-statistics representation consistently
outperforms 30--200 dimensional spectral baselines. Improvements range
from +3.0 to +11.1 percentage points in overall accuracy and +6.3 to
+16.1 points in average accuracy, with particularly strong gains for
minority classes.

Performance saturates at a small, scene-dependent neighborhood scale
(5\(\times\)5--7\(\times\)7), suggesting the existence of a minimal sufficient spatial
radius beyond which larger windows provide no consistent benefit. In
contrast, uncertainty-driven multimodal gating strategies fail to
improve performance and, in some cases, substantially degrade it, with
gate confidence exhibiting slight negative correlation with
classification correctness.

These results indicate that local second-order spatial statistics
constitute a dominant and efficient representational mechanism for
standard hyperspectral benchmarks. Importantly, this representation
achieves superior performance with substantial dimensionality reduction
(76--95\% fewer features) and minimal computational overhead.

Future work should evaluate this principle on additional multimodal
datasets, explore adaptive or anisotropic neighborhood selection, and
analyze the interaction between explicit neighborhood statistics and
learned convolutional representations.

\end{document}
